\subsection{End-to-end simulation of a stochastic Transformer}
\label{sec:full_transformer}

We next integrated the stochastic operators into a complete four-block
encoder-only Transformer. The network included learned input and positional
embeddings, two attention heads per block, output projections, residual
connections, pre-layer normalization, feed-forward networks with GELU
activations, a final normalization layer and a classification head. We used
a controlled associative-retrieval task: 15 tokens presented distinct keys
with independently assigned four-class values, and a final query token
requested the value associated with one stored key. Three independently
trained floating-point models reached 100\% accuracy on the 128 held-out
sequences used for stochastic evaluation.

All main linear stages and attention products were evaluated using explicit
sign--magnitude bitstreams. GELU used an eight-state probability gate, whereas
layer normalization combined stochastic variance estimation with a four-state
reciprocal-square-root gate. We examined direct transfer of the softmax
operator from Section~\ref{sec:softmax_approximation} and a simple
activation-calibration procedure. Calibration updated the same-size
HI-SMURF table using reference attention scores from separate input sequences,
with its encoding range selected on an independent calibration-validation set.
The Transformer weights remained fixed, and the HI-SMURF operator retained
16 states, TT rank 8 and 16,512 parameters. Calibration used reference
operator outputs without task labels.

Calibration reduced errors throughout the network
(Fig.~\ref{fig:transformer_depth}). At $L=256$, the relative error in the
query representation after the fourth block decreased from 0.736
under direct transfer to 0.122 after calibration. We additionally
measured the error in each block's residual update to assess the learned
transformation separately from the identity connection. At the classifier
output, relative centred-logit error decreased from
0.811 $\pm$ 0.081 to 0.079 $\pm$ 0.021,
and retrieval accuracy increased from 71.4\% to 99.9\%
(Fig.~\ref{fig:transformer_performance}). Values denote means across three
trained models after averaging three stream realizations per model; reported
uncertainty is the sample standard deviation across the three model means.

Longer streams improved numerical fidelity over the tested range, with
diminishing gains within the calibrated architecture.
Increasing $L$ from 64 to 1,024 reduced relative centred-logit error from
0.230 to 0.081, compared with 0.075 for analytic
evaluation of the calibrated finite-state operators
(Fig.~\ref{fig:transformer_performance}b). Error did not decrease monotonically:
the mean at $L=1024$ was slightly above that at $L=256$ (0.079),
indicating that longer streams did not remove the remaining approximation
error. Corresponding retrieval
accuracies were 98.0\% and 99.7\%
(Fig.~\ref{fig:transformer_performance}a). The continuous-output metric
therefore resolves approximation errors that need not change the predicted
class. Stream length excludes 256 burn-in cycles for each multivariate
softmax FSM ensemble and 64 cycles for each scalar FSM stage.

Progressively extending stochastic computation from the attention cores to
all learned linear stages and then to normalization and GELU yielded relative
centred-logit errors of 0.093, 0.087 and 0.079, respectively,
at $L=256$ (Fig.~\ref{fig:transformer_ablation}a). These end-to-end controls
contain interacting approximation errors; their ordering does not establish
that adding stochastic stages improves intrinsic accuracy. With the calibrated
operator fixed, switching from independent arithmetic streams to Sobol
streams changed this error from 0.249 to 0.079
(Fig.~\ref{fig:transformer_ablation}b). These comparisons evaluate the
complete encoder graph on a controlled synthetic task. TT contractions,
centering, residual and bias accumulation, static normalization-affine gains,
range scaling and conversion between stochastic stages remained numerical
operations; hardware area, energy and latency were not measured.
